% !TeX root = ../main.tex
\section{Seno, coseno e tangente}
% Fonte: Lezione 6, pp.~10 e 21

Gli angoli sono misurati in \emph{radianti}: un giro completo misura $2\pi$. Sulla circonferenza di centro l'origine e raggio $1$, il punto individuato dall'angolo orientato $x$ ha coordinate
\[
P=(\cos x,\sin x).
\]
Il coseno è dunque l'ascissa e il seno l'ordinata del punto.

\begin{schema}
\begin{tikzpicture}[scale=2.1]
  \draw[asse] (-1.25,0) -- (1.55,0) node[right] {$u$};
  \draw[asse] (0,-1.15) -- (0,1.25) node[above] {$v$};
  \draw[black!55] (0,0) circle (1);
  \draw[orange!80!black,thick] (0,0) -- (1,{tan(35)});
  \draw[coloreschema,thick] (0,0) -- ({cos(35)},{sin(35)});
  \draw[black!45] (1,-.15) -- (1,1.05);
  \draw[dotted] ({cos(35)},0) -- ({cos(35)},{sin(35)})
    -- (0,{sin(35)});
  \draw[->] (.3,0) arc[start angle=0,end angle=35,radius=.3];
  \node at (.41,.12) {$x$};
  \fill[coloreschema] ({cos(35)},{sin(35)}) circle (.8pt)
    node[above left] {$P$};
  \fill[orange!80!black] (1,{tan(35)}) circle (.8pt);
  \node[below] at ({cos(35)},0) {$\cos x$};
  \node[left] at (0,{sin(35)}) {$\sin x$};
  \node[right] at (1.03,.35) {$\tan x$};
  \node[below left] at (0,0) {$O$};
  \node[below right] at (1,-.12) {$1$};
\end{tikzpicture}

Nel disegno $0<x<\pi/2$. La retta di pendenza $\tan x$ incontra la tangente verticale $u=1$ all'altezza $\tan x$.
\end{schema}

\textbf{Dominio, immagine e periodicità.} Seno e coseno sono definiti su tutto $\R$, hanno immagine $[-1,1]$ e periodo principale $2\pi$. La periodicità esprime il ritorno allo stesso punto dopo un giro completo.

\textbf{Parità.} Riflettendo il punto della circonferenza rispetto all'asse orizzontale si ottiene
\[
\sin(-x)=-\sin x,\qquad \cos(-x)=\cos x:
\]
il seno è dispari, il coseno è pari.

\textbf{Identità fondamentale.} Dall'equazione della circonferenza unitaria, o dal teorema di Pitagora, segue l'identità fondamentale
\[
\boxed{\sin^2x+\cos^2x=1}\qquad\forall x\in\R.
\]


\textbf{Estremi.} I valori $\pm1$ sono raggiunti, quindi sono anche minimo e massimo delle due funzioni.

\textbf{Tangente.} La \emph{tangente}, indicata anche con $\operatorname{tg}$, è
\[
\tan x=\frac{\sin x}{\cos x},\qquad
\operatorname{dom}(\tan)=\R\setminus
\left\{\frac\pi2+k\pi:k\in\Z\right\}.
\]
Sono esclusi esattamente gli zeri del denominatore.

\textbf{Parità e periodo.} La tangente è dispari e ha periodo principale $\pi$: seno e coseno cambiano entrambi segno dopo mezzo giro, lasciando invariato il loro rapporto.

\textbf{Immagine e monotonia.} La sua immagine è $\R$, quindi non è limitata. Su ogni intervallo $(-\pi/2+k\pi,\pi/2+k\pi)$ è strettamente crescente e assume tutti i valori reali; non è monotona sull'intero dominio, che contiene più rami distinti.

\section{Restrizioni invertibili e arcoseno}
% Fonte: Lezione 6, p.~22

Seno e coseno sono periodici e perciò non sono iniettivi su $\R$. Per costruirne le inverse si sceglie un intervallo in cui ciascuna funzione sia strettamente monotona e assuma l'intera immagine $[-1,1]$.

\textbf{Ramo invertibile del seno.} Il seno è strettamente crescente da $-1$ a $1$ su $[-\pi/2,\pi/2]$. Sulla semicirconferenza destra, infatti, l'ordinata del punto aumenta al crescere dell'angolo in tale intervallo.

\textbf{Arcoseno.} L'inversa della restrizione è l'\emph{arcoseno}:
\[
\arcsin\colon[-1,1]\longrightarrow
\left[-\frac\pi2,\frac\pi2\right],
\qquad \arcsin y=\text{l'unico }x\in[-\pi/2,\pi/2]\text{ con }\sin x=y.
\]
È strettamente crescente e dispari; il suo grafico è quello del ramo scelto del seno riflesso rispetto a $y=x$.

\begin{schema}
\begin{tikzpicture}[x=.85cm,y=1.15cm]
  \draw[asse] (-3.6,0) -- (3.6,0) node[right] {$x$};
  \draw[asse] (0,-1.4) -- (0,1.5) node[above] {$y$};
  \draw[black!40,domain=-3.4:3.4,samples=100] plot (\x,{sin(\x r)});
  \draw[intervallo,domain=-pi/2:pi/2,samples=70] plot (\x,{sin(\x r)});
  \draw[dotted] (-1.571,0) -- (-1.571,-1) -- (0,-1);
  \draw[dotted] (1.571,0) -- (1.571,1) -- (0,1);
  \fill[coloreschema] (-1.571,-1) circle (1.4pt);
  \fill[coloreschema] (1.571,1) circle (1.4pt);
  \node[above] at (-1.571,0) {$-\frac\pi2$};
  \node[below] at (1.571,0) {$\frac\pi2$};
  \node[above right] at (0,1) {$1$};
  \node[below right] at (0,-1) {$-1$};
  \node[below] at (0,-1.8) {Seno; in colore il ramo invertibile};
\end{tikzpicture}
\hspace{.6cm}
\begin{tikzpicture}[x=1.75cm,y=1cm]
  \draw[asse] (-1.35,0) -- (1.4,0) node[right] {$x$};
  \draw[asse] (0,-1.9) -- (0,1.95) node[above] {$y$};
  \draw[intervallo,domain=-pi/2:pi/2,samples=100] plot ({sin(\x r)},\x);
  \draw[dotted] (-1,0) -- (-1,-1.571) -- (0,-1.571);
  \draw[dotted] (1,0) -- (1,1.571) -- (0,1.571);
  \fill[coloreschema] (-1,-1.571) circle (1.4pt);
  \fill[coloreschema] (1,1.571) circle (1.4pt);
  \node[above] at (-1,0) {$-1$};
  \node[below] at (1,0) {$1$};
  \node[left] at (0,1.571) {$\frac\pi2$};
  \node[right] at (0,-1.571) {$-\frac\pi2$};
  \node[below] at (0,-2.07) {Arcoseno};
\end{tikzpicture}
\end{schema}

\textbf{Identità di inversione e domini.} Si ha
\[
\sin(\arcsin y)=y\quad(y\in[-1,1]),\qquad
\arcsin(\sin x)=x\quad(x\in[-\pi/2,\pi/2]).
\]
\begin{esempio}
La seconda identità non vale per ogni reale: $\arcsin(\sin\pi)=\arcsin0=0\ne\pi$.
\end{esempio}

\section{Arcocoseno}
% Fonte: Lezione 6, p.~23

\textbf{Ramo invertibile del coseno.} Il coseno è strettamente decrescente da $1$ a $-1$ su $[0,\pi]$: percorrendo la semicirconferenza superiore in senso antiorario, l'ascissa diminuisce.

\textbf{Arcocoseno.} Si definisce l'\emph{arcocoseno} come l'inversa di questa restrizione:
\[
\arccos\colon[-1,1]\longrightarrow[0,\pi],\qquad
\arccos y=\text{l'unico }x\in[0,\pi]\text{ con }\cos x=y.
\]
\textbf{Monotonia e identità di inversione.} È strettamente decrescente e soddisfa
\[
\cos(\arccos y)=y\quad(y\in[-1,1]),\qquad
\arccos(\cos x)=x\quad(x\in[0,\pi]).
\]
\textbf{Valori notevoli.} Si ha $\arccos1=0$, $\arccos0=\pi/2$ e $\arccos(-1)=\pi$.

\begin{schema}
\begin{tikzpicture}[x=.85cm,y=1.15cm]
  \draw[asse] (-1.9,0) -- (4.5,0) node[right] {$x$};
  \draw[asse] (0,-1.4) -- (0,1.5) node[above] {$y$};
  \draw[black!40,domain=-1.7:4.3,samples=100] plot (\x,{cos(\x r)});
  \draw[intervallo,domain=0:pi,samples=70] plot (\x,{cos(\x r)});
  \draw[dotted] (3.142,0) -- (3.142,-1) -- (0,-1);
  \fill[coloreschema] (0,1) circle (1.4pt);
  \fill[coloreschema] (3.142,-1) circle (1.4pt);
  \node[above] at (3.142,0) {$\pi$};
  \node[left] at (0,1) {$1$};
  \node[left] at (0,-1) {$-1$};
  \node[below] at (1.2,-1.8) {Coseno; in colore il ramo invertibile};
\end{tikzpicture}
\hspace{.6cm}
\begin{tikzpicture}[x=1.75cm,y=1cm]
  \draw[asse] (-1.35,0) -- (1.4,0) node[right] {$x$};
  \draw[asse] (0,-.3) -- (0,3.55) node[above] {$y$};
  \draw[intervallo,domain=0:pi,samples=100] plot ({cos(\x r)},\x);
  \draw[dotted] (-1,0) -- (-1,3.142) -- (0,3.142);
  \fill[coloreschema] (-1,3.142) circle (1.4pt);
  \fill[coloreschema] (1,0) circle (1.4pt);
  \node[below] at (-1,0) {$-1$};
  \node[below] at (1,0) {$1$};
  \node[right] at (0,3.142) {$\pi$};
  \node[left] at (0,1.571) {$\frac\pi2$};
  \node[below] at (0,-.92) {Arcocoseno};
\end{tikzpicture}
\end{schema}

\Needspace{31\baselineskip}
\section{Arcotangente}
% Fonte: Lezione 6, p.~24

\textbf{Ramo invertibile della tangente.} La restrizione $\tan\colon(-\pi/2,\pi/2)\to\R$ è biettiva. Geometricamente, mentre il punto percorre la semicirconferenza destra, la retta che passa per l'origine incontra la verticale $u=1$ a un'altezza che cresce da valori negativi arbitrariamente grandi in modulo a valori positivi arbitrariamente grandi. Ogni altezza reale è raggiunta una sola volta.

\textbf{Arcotangente.} Si definisce l'\emph{arcotangente}:
\[
\arctan\colon\R\longrightarrow\left(-\frac\pi2,\frac\pi2\right),
\qquad \arctan y=\text{l'unico }x\in(-\pi/2,\pi/2)\text{ con }\tan x=y.
\]
\textbf{Proprietà e identità di inversione.} È strettamente crescente, dispari e limitata. Vale
\[
\tan(\arctan y)=y\quad(y\in\R),\qquad
\arctan(\tan x)=x\quad(x\in(-\pi/2,\pi/2)).
\]
\textbf{Estremi non raggiunti.} Il suo supremo è $\pi/2$ e il suo infimo è $-\pi/2$, ma non ha né massimo né minimo: gli estremi dell'intervallo immagine sono esclusi. Le rette $y=\pm\pi/2$ indicano le altezze a cui il grafico si avvicina per argomenti rispettivamente molto grandi positivi e negativi.

\begin{schema}
\begin{tikzpicture}[x=.85cm,y=.65cm]
  \draw[asse] (-2.3,0) -- (4.9,0) node[right] {$x$};
  \draw[asse] (0,-3) -- (0,3.2) node[above] {$y$};
  \foreach \a in {-1.571,1.571,4.712}
    \draw[dashed,black!40] (\a,-2.8) -- (\a,2.8);
  \draw[intervallo,domain=-1.22:1.22,samples=90] plot (\x,{tan(\x r)});
  \draw[black!45,domain=1.93:4.36,samples=90] plot (\x,{tan(\x r)});
  \node[below left] at (-1.571,0) {$-\frac\pi2$};
  \node[below right] at (1.571,0) {$\frac\pi2$};
  \node[below] at (4.712,0) {$\frac{3\pi}2$};
  \node[below] at (1.2,-3.5) {Tangente; in colore il ramo invertibile};
\end{tikzpicture}
\hspace{.6cm}
\begin{tikzpicture}[x=.65cm,y=1.1cm]
  \draw[asse] (-3.4,0) -- (3.5,0) node[right] {$x$};
  \draw[asse] (0,-1.95) -- (0,1.95) node[above] {$y$};
  \draw[dashed,black!45] (-3.2,1.571) -- (3.2,1.571);
  \draw[dashed,black!45] (-3.2,-1.571) -- (3.2,-1.571);
  \draw[intervallo,domain=-3.2:3.2,samples=100] plot (\x,{atan(\x)*pi/180});
  \node[above right] at (0,1.571) {$\frac\pi2$};
  \node[below right] at (0,-1.571) {$-\frac\pi2$};
  \node[below] at (0,-2.17) {Arcotangente};
\end{tikzpicture}
\end{schema}
